% Part: computability
% Chapter: recursive-functions
% Section: primitive-recursion

\documentclass[../../../include/open-logic-section]{subfiles}

\begin{document}

\olfileid{cmp}{rec}{pre}
\olsection{આદિમ પુનરાવર્તન}

પ્રાકૃતિક સંખ્યાઓની એક ખાસિયત એ છે કે $0$થી શરૂ કરીને ઉત્તરગામી
ક્રિયા~$+1$ સાન્ત વખત કરતાં દરેક પ્રાકૃતિક સંખ્યા સુધી પહોંચી શકાય
છે---દરેક પ્રાકૃતિક સંખ્યા કાં તો~$0$ છે અથવા~$0$ના ઉત્તરગામીના
\dots{} ઉત્તરગામી છે. આ હકીકતના આધારે $h\colon\Nat \to \Nat$
વિધેય આપવાની એક રીત આવી છે: (a)~આર્ગ્યુમેન્ટ~$0$ માટે $h$નું મૂલ્ય
જણાવો, અને (b)~$h(x)$નું મૂલ્ય જાણ્યા પછી $h(x+1)$નું મૂલ્ય કેવી રીતે
ગણવું તે પણ જણાવો. (a) સીધું $h(0)$ આપે છે, તેથી $h$નું મૂલ્ય~$0$ માટે
નક્કી થાય છે. હવે (b)ની સૂચના $x=0$ માટે વાપરીને~$h(0)$માંથી
$h(1) = h(0+1)$ ગણી શકાય. એ જ સૂચના $x=1$ માટે વાપરીને~$h(1)$માંથી
$h(2) = h(1+1)$ ગણાય, અને આમ આગળ વધીએ. દરેક પ્રાકૃતિક સંખ્યા~$x$
માટે અંતે એવો તબક્કો આવે છે જેમાં $h(x)$માંથી $h(x+1)$ નક્કી કરીએ
છીએ. તેથી દરેક $x \in \Nat$ માટે $h(x)$ વ્યાખ્યાયિત થાય છે.

\sourcecorrection{OLCMP-001}{સ્થિર અંગ્રેજી મૂળના આ વાક્યમાં
``$h(x+1)$માંથી $h(x)$ નક્કી કરીએ'' એવી દિશા ઊલટી લખાઈ છે.
અગાઉ આપેલી (b) સૂચના મુજબ $h(x)$માંથી $h(x+1)$ નક્કી થાય છે;
ગુજરાતી વાક્યમાં આ દિશા સુધારી છે.}

દાખલા તરીકે, ધારો કે નીચેના બે સમીકરણોથી $h\colon \Nat \to \Nat$
આપ્યું છે:
\begin{align*}
h(0) & =  1\\
h(x+1) & =  2 \cdot h(x)
\end{align*}
જો ગુણાકાર કેવી રીતે કરવો તે પહેલેથી જાણતા હોઈએ, તો આ સમીકરણો ઉપરની
(a) અને~(b) માટે જરૂરી માહિતી આપે છે. બીજું સમીકરણ વારંવાર વાપરતાં
મળે છે કે
\begin{align*}
  h(1) & = 2\cdot h(0) = 2,\\
  h(2) & = 2\cdot h(1) = 2\cdot 2,\\
  h(3) & = 2 \cdot h(2) = 2\cdot 2 \cdot 2,\\
  & \vdots
\end{align*}
આથી આપણે આપેલું વિધેય~$h$ એટલે $h(x) = 2^x$ છે.

પ્રાકૃતિક સંખ્યાઓની ઉપર જણાવેલી ખાસિયત ખાતરી આપે છે કે આ બે શરતો
પૂરી કરતું $h$ વિધેય એક જ છે. આવા બે સમીકરણોની જોડીને $h$ની
\emph{આદિમ પુનરાવર્તન દ્વારા વ્યાખ્યા} કહે છે. આ રીતે $h$ને
નક્કી કરીએ છીએ. તેને ``પુનરાવર્તી''
કહેવાય છે, કારણ કે વ્યાખ્યામાં---ખાસ કરીને બીજા સમીકરણની જમણી
બાજુએ---$h$ પોતે આવે છે. તેને ``આદિમ'' કહેવાય છે, કારણ કે
$h(x+1)$ નક્કી કરતી વખતે માત્ર તરત પહેલાંનું મૂલ્ય~$h(x)$ વાપરીએ
છીએ. $\Nat$ પર વિધેયને પુનરાવર્તી રીતે વ્યાખ્યાયિત કરવાની આ સૌથી
સરળ રીત છે.

સરવાળો અને ગુણાકાર જેવા વધુ મૂળભૂત વિધેયો પણ આદિમ પુનરાવર્તનથી
વ્યાખ્યાયિત કરી શકાય છે. જોકે અહીં સંબંધિત વિધેયોને $2$ આર્ગ્યુમેન્ટ
છે. આપણે એક આર્ગ્યુમેન્ટનું સ્થાન નિશ્ચિત રાખીએ અને બીજા પર
પુનરાવર્તન કરીએ. દાખલા તરીકે, $\Add(x, y)$ વ્યાખ્યાયિત કરવા~$x$
નિશ્ચિત રાખીને પહેલાં $y=0$ માટે મૂલ્ય આપીએ અને પછી~$y$ માટેના
મૂલ્ય પરથી $y+1$ માટેનું મૂલ્ય આપીએ. $x$ નિશ્ચિત છે, તેથી તે
વ્યાખ્યાકારક સમીકરણોની ડાબી અને જમણી બંને બાજુ આવે છે.
\begin{align*}
\Add(x,0) & =  x\\
\Add(x,y+1) & =  \Add(x,y)+1
\end{align*}
આ સમીકરણો દરેક $x$ \emph{અને}~$y$ માટે $\Add$નું મૂલ્ય આપે છે.
દાખલા તરીકે $\Add(2,3)$ શોધવા $x = 2$ રાખીને પહેલાં પ્રથમ સમીકરણથી
$\Add(2,0) = 2$ મળે છે. પછી બીજા સમીકરણને ક્રમે વાપરતાં
$\Add(2,1) = 2 + 1 = 3$, $\Add(2, 2) = 3 + 1 = 4$ અને $\Add(2,
3) = 4 + 1 = 5$ મળે છે.

$\Add$ની વ્યાખ્યાના બીજા સમીકરણની જમણી બાજુએ આપણે $+$ વાપર્યું,
પણ તે માત્ર~$1$ ઉમેરવા માટે છે. બીજા શબ્દોમાં, ઉત્તરગામી વિધેય
$\Succ(z) = z+1$ને અગાઉના મૂલ્ય $\Add(x,y)$ પર લાગુ કરીને
$\Add(x,y+1)$ આપ્યું છે. તેથી આ પુનરાવર્તી વ્યાખ્યાને અગાઉના
મૂલ્ય પર લાગુ થતા એક વિધેય દ્વારા આપવામાં આવેલી માની શકાય. છતાં
એ વિધેય માત્ર અગાઉના મૂલ્ય પર જ નહીં, પણ $x$ અને~$y$ પર પણ આધાર
રાખે એવી છૂટ આપવાથી કોઈ અડચણ નથી---અને ક્યારેક એવી છૂટ જરૂરી છે.
આ જુઓ:
\begin{align*}
  \Mult(x,0) & =  0 \\
  \Mult(x,y+1) & =  \Add(\Mult(x,y),x)
\end{align*}
અહીં અગાઉના મૂલ્ય $\Mult(x,y)$ અને પ્રથમ આર્ગ્યુમેન્ટ~$x$ પર
$\Add$ લાગુ કરીને $\Mult$ની આદિમ પુનરાવર્તી વ્યાખ્યા આપી છે.
આ વ્યાખ્યા દરેક $x$ અને~$y$ માટે $\Mult(x,y)$ નક્કી કરે છે.
દાખલા તરીકે $\Mult(2,3)$, ક્રમે $\Mult(2,0)$, $\Mult(2,1)$,
$\Mult(2,2)$ અને~$\Mult(2,3)$ ગણીને મળે છે:
\begin{align*}
  \Mult(2,0) & = 0\\
  \Mult(2,1) & = \Mult(2,0+1) =
  \Add(\Mult(2,0), 2) = \Add(0, 2) = 2\\
  \Mult(2,2) & = \Mult(2,1+1) =
  \Add(\Mult(2,1), 2) = \Add(2, 2) = 4\\
  \Mult(2,3) & = \Mult(2,2+1) =
  \Add(\Mult(2,2), 2) = \Add(4, 2) = 6
\end{align*}

સામાન્ય માળખું આ પ્રમાણે છે. $h(x_0, \dots, x_{k-1}, y)$ની આદિમ
પુનરાવર્તી વ્યાખ્યા આપવા બે સમીકરણો જોઈએ. પહેલું સમીકરણ~$h$નો
ઉલ્લેખ કર્યા વિના $h(x_0, \dots, x_{k-1}, 0)$નું મૂલ્ય આપે છે.
બીજું સમીકરણ $h(x_0, \dots, x_{k-1}, y)$, અન્ય આર્ગ્યુમેન્ટ
$x_0$, \dots,~$x_{k-1}$ અને~$y$ના આધારે
$h(x_0, \dots, x_{k-1}, y+1)$નું મૂલ્ય આપે છે. બીજા સમીકરણમાં
$h$નું માત્ર તરત પહેલાંનું મૂલ્ય જ વાપરી શકાય. જો બંને સમીકરણોની
જમણી બાજુએ આવતી ક્રિયાઓને અનુક્રમે $f$ અને~$g$ વિધેયો માનીએ,
તો આદિમ પુનરાવર્તનથી નવું વિધેય~$h$ વ્યાખ્યાયિત કરવાની સામાન્ય
પદ્ધતિ આ છે:
\begin{align*}
  h(x_0, \dots, x_{k-1}, 0) & = f(x_0, \dots, x_{k-1})\\
  h(x_0, \dots, x_{k-1}, y+1) & =
  g(x_0, \dots, x_{k-1}, y, h(x_0, \dots, x_{k-1}, y))
\end{align*}
$\Add$ માટે $k=1$, $f(x_0) = x_0$ (તદેવ વિધેય) અને
$g(x_0, y, z) = z + 1$ છે (પોતાના ત્રીજા આર્ગ્યુમેન્ટનો ઉત્તરગામી
આપતું $3$ આર્ગ્યુમેન્ટવાળું વિધેય):
\begin{align*}
  \Add(x_0, 0) & = f(x_0) = x_0\\
  \Add(x_0, y+1) & = g(x_0, y, \Add(x_0, y)) =
  \Succ(\Add(x_0, y))
\end{align*}
$\Mult$ માટે $f(x_0) = 0$ (હંમેશાં~$0$ આપતું અચળ વિધેય) અને
$g(x_0, y, z) = \Add(z,x_0)$ છે (પોતાના છેલ્લા અને પ્રથમ
આર્ગ્યુમેન્ટનો સરવાળો આપતું $3$ આર્ગ્યુમેન્ટવાળું વિધેય):
\begin{align*}
  \Mult(x_0,0) & =  f(x_0) = 0 \\
  \Mult(x_0,y+1) & =  g(x_0, y, \Mult(x_0,y)) =
  \Add(\Mult(x_0,y), x_0)
\end{align*}

\end{document}
